% ============================================================================
% CHAPITRE 3 — SECTION 3.5 : CONCLUSION
% Fichier : partie3/conclusion_validation.tex
% Aligné sur la structure visée : 3.5 Conclusion du Chapitre 3
% ============================================================================

\section{CONCLUSION DU CHAPITRE~\ref{chap:experimentation}}
\label{sec:conclusion_experimentation}
% ============================================================================
\begin{sloppypar}

Ce chapitre, qui correspond \`a l'\'etape \textbf{X} de la progression
F$\to$B$\to$E$\to$S$\to$C$\to$\textbf{X}, a trait\'e une question que les bornes du
chapitre~\ref{chap:cadre_formel} laissaient enti\`ere, celle de la valeur
d'in\'egalit\'es dont les param\`etres ne sont pas mesur\'es. Il en r\'esulte trois acquis
assortis d'une r\'eserve.

Le premier acquis est m\'ethodologique. Les six param\`etres
institutionnels deviennent des grandeurs mesurables. Chacun dispose
d\'esormais d'une d\'efinition op\'erationnelle, d'un estimateur explicite,
d'un plan d'\'echantillonnage chiffr\'e et de crit\`eres de fiabilit\'e
emprunt\'es aux standards des sciences du comportement, tels que l'accord
inter-\'evaluateurs, la fid\'elit\'e test-retest et les seuils ICC et $\kappa$. La
cryptographie institutionnelle dispose ainsi de son protocole de mesure avant
d'avoir ses mesures. Cet ordre s'impose puisqu'un r\'esultat n'est
falsifiable que si la proc\'edure de falsification est sp\'ecifi\'ee \`a
l'avance.

Le deuxi\`eme acquis est une calibration. La d\'erivation descendante
(section~\ref{sec:derivation_opposabilite}) part des espaces de verdicts, tandis
que la d\'ecomposition des capacit\'es d'extraction
(section~\ref{sec:estimation_params_classiques}) part des ressources des
acteurs. Les deux voies aboutissent \`a des ordres de grandeur compatibles, r\'eunis dans une
table unique de r\'ef\'erence (tableau~\ref{tab:parametres_reference}). Cette
compatibilit\'e ne vaut pas validation. Les deux voies ne sont pas
ind\'ependantes, puisque les valeurs de $\Delta$ ont d'abord \'et\'e fix\'ees par
inversion de la borne. Les syst\`emes d\'eploy\'es ne satisfont pas non plus toutes les
hypoth\`eses du th\'eor\`eme (Remarque~\ref{rem:portee_cas_etude}). Cette compatibilit\'e
montre seulement que les grandeurs retenues ne se contredisent pas.

Le troisi\`eme acquis, qui est n\'egatif, se r\'ev\`ele le plus instructif. Sur les
standards post-quantiques, les produits CLO$^q$ illustratifs d\'epassent la borne
(tableau~\ref{tab:decomposition_borne_pq})~; il en va de m\^eme de l'\'etude de cas
ML-DSA. Ces d\'epassements ne r\'efutent pas l'extension quantique, dont les
syst\`emes consid\'er\'es ne sont pas des instances~; ils montrent que les
param\`etres $\Delta_q$ retenus sont des postulats qui constituent la premi\`ere
cible du protocole de la section~\ref{subsec:protocole_delta_q}.

La r\'eserve, enfin, a d\'ej\`a \'et\'e \'enonc\'ee \`a la
section~\ref{subsubsec:statut_deploiement} et ne sera pas att\'enu\'ee ici.
Le programme exp\'erimental sp\'ecifi\'e n'a pas \'et\'e d\'eploy\'e.
Tant qu'il ne l'aura pas \'et\'e, les Trilemmes CRO conservent le
statut que leur \'enonc\'e leur conf\`ere, celui de r\'esultats math\'ematiquement
\'etablis sous hypoth\`eses, dont les param\`etres institutionnels restent \`a
mesurer. Le
chapitre~\ref{chap:resultats_discussions} interpr\`ete ces bornes et en tire
les implications~; il le fait dans les limites que le pr\'esent
chapitre~\ref{chap:experimentation} vient de tracer.

\end{sloppypar}